\chapter{Rules of correspondence with ERA5}
\label{cap:regole}

Two rules decide what the archive publishes and where, and they apply to all 90 variables.

\section{Publication rule}
\label{sec:pubblicazione}

A field is published under an ERA5 name only if the WRF quantity is the one that the ERA5
parameter defines, or if the difference can be stated in one line. A field that meets neither
condition is left out of the archive, even when the model computes it. There are two distinct
reasons for an exclusion. For some WRF fields no ECMWF parameter
describes the same quantity, for others the parameter exists but ERA5 does not publish it
(\cref{sec:non-pubblicati}).

Seven fields, in six rows of \cref{tab:approssimazioni}, are published with a declared difference
from their ERA5 namesake.

\begin{table}[htbp]
  \centering
  \small
  \caption{Fields published with a declared difference from the ERA5 parameter.}
  \label{tab:approssimazioni}
  \begin{tabularx}{\textwidth}{@{}lXXl@{}}
    \toprule
    Field & ERA5 parameter & CHAPTER field & Section \\
    \midrule
    \campo{slor} & slope of the sub-grid orography & slope of the terrain resolved at \SI{3}{km} & \ref{sec:orografia} \\
    \campo{10fg} & parameterized gust & hourly maximum of the resolved wind at \SI{10}{m} & \ref{cap:raffica} \\
    \campo{iews}, \campo{inss} & components of the IFS turbulent surface stress & $\rho u_*^2$, with $\rho$ the air density and $u_*$ the friction velocity, projected onto the direction of the \SI{10}{m} wind & \ref{sec:rugosita} \\
    \campo{fsr} & effective roughness length, including the parameterized orographic drag & local roughness length \campo{ZNT}, table value of the dominant category modified by snow & \ref{sec:rugosita} \\
    \campo{lsm} & permanent land--sea mask & hourly WRF mask, which counts sea ice as land & \ref{sec:statici-orari} \\
    \campo{tcw} & total column water, five species & sum of six species, graupel included & \ref{sec:colonne} \\
    \bottomrule
  \end{tabularx}
\end{table}

\section{Masking rule}
\label{sec:mascheratura}

A point is missing where the ERA5 parameter is not defined. A point where the model provides no
value, or where the value looks wrong, is not missing for that reason. Missing points are written
as a GRIB bitmap, never as zeros, so a reader of the archive can always tell a point where the
parameter is not defined from one where it is zero. \Cref{tab:maschere} lists the masked fields.

\begin{table}[htbp]
  \centering
  \small
  \caption{Fields with missing points and the domain on which they are defined.}
  \label{tab:maschere}
  \begin{tabularx}{\textwidth}{@{}lXl@{}}
    \toprule
    Field & Defined on & Section \\
    \midrule
    \campo{sst}, \campo{ci} & sea cells & \ref{sec:mare} \\
    \campo{dl} & cells with a lake cover greater than zero (\num{32000}) & \ref{sec:vegetazione} \\
    \campo{swvl1-4}, \campo{stl1-4}, \campo{slt} & permanent land (\num{1304109} cells) & \ref{sec:suolo}, \ref{sec:vegetazione} \\
    \campo{tvl} & all cells except the \num{124672} where shrubland dominates & \ref{sec:vegetazione} \\
    \campo{mucin} & columns with CAPE of at least \SI{100}{J/kg} & \ref{sec:instabilita} \\
    \campo{tsn} & cells with \campo{SNOWC} above \num{0.9} & \ref{sec:neve} \\
    \campo{fal} & points with incoming solar radiation at the surface above \SI{50}{W/m^2} & \ref{sec:albedo} \\
    \bottomrule
  \end{tabularx}
\end{table}

The lake cover \campo{cl} is defined over the sea as well, where it is zero, because a sea cell
has no lake cover. Four fields are masked for a different reason: at the missing points the WRF
quantity has no value that corresponds to the ERA5 parameter. \campo{tvl} is missing over
shrubland, because the ERA5 type table separates evergreen and deciduous shrubs and the MODIS
classification carries no shrub phenology (\cref{sec:vegetazione}). \campo{tsn} is missing where
snow does not cover the cell, because there the model field is a soil temperature
(\cref{sec:neve}). \campo{fal} is missing where the incoming solar radiation at the surface does not
exceed \SI{50}{W/m^2}; without incoming radiation the ratio of reflected to incoming radiation is
not defined (\cref{sec:albedo}). \campo{mucin}
is missing where CAPE is below \SI{100}{J/kg}, because there the routine returns no inhibition
(\cref{sec:instabilita}).

\subsection{Sea and inland water}
\label{sec:mare}

Masking \campo{sst} and \campo{ci} requires separating the sea from inland water, and the lake
class of the static file is not sufficient for this. Class 21 (lake) of the MODIS-IGBP
classification includes the Black Sea, the Sea of Azov and part of the Baltic, which together make
up \SI{89}{\percent} of the \num{55881} cells that the static file classifies entirely as lake.
Masking on that class would remove \campo{sst} from the Black Sea.

The criterion used is the size of the connected component, i.e.\ of the set of water cells
connected through shared edges. In the domain the largest components have \num{844431} cells (Atlantic,
Mediterranean, Baltic and North Sea), \num{49841} (Black Sea and Sea of Azov) and \num{8583} (Red
Sea). The next is Lake Vänern, with \num{1225} cells, the largest lake in the domain. Components of
at least \num{2000} cells are sea, for a total of \num{902855} sea cells. Any threshold between
\num{1226} and \num{8583} cells gives the same set.

\section{Static fields and hourly fields}
\label{sec:statici-orari}

A static field comes from the static file and is therefore identical in every hourly file. The
static fields are \campo{cl}, \campo{cvh}, \campo{cvl}, \campo{dl}, \campo{slt}, \campo{tvh} and
\campo{tvl}. The soil columns \campo{swvl1-4} and \campo{stl1-4} change every hour, but their mask
is the permanent land of the static file.

The land--sea mask \campo{lsm}, by contrast, is the hourly WRF mask (\campo{LANDMASK}), which
reclassifies as land every water cell covered by sea ice. The mask therefore grows in winter: in a
file of 20 March 2024 the reclassified cells number \num{2469}, exactly the cells with sea ice.
Sea ice \campo{ci} is 0 or 1, because the simulation does not use fractional sea ice
(\campo{FRACTIONAL_SEAICE=0}). The permanent land mask, in the ERA5 sense, is the set of points
with $\campo{lsm} = 1$ and $\campo{ci} = 0$, counting the missing points of \campo{ci} as zero.

The background albedo \campo{al} also changes in time, with sea ice and at the transition between
the summer and winter seasons of the land-use table (\cref{sec:albedo}). For \campo{lsm} and
\campo{al} must be read from the hourly file of the validity time, not from an arbitrary file of
the archive.
